\documentclass[11pt]{amsart}
\usepackage[utf8]{inputenc}
\usepackage[T1]{fontenc}
\usepackage{lmodern}
\usepackage[margin=1.05in]{geometry}
\usepackage{amsmath,amssymb,amsthm,mathtools}
\usepackage{booktabs}
\usepackage{array}
\usepackage{enumitem}
\usepackage[hidelinks]{hyperref}
\usepackage{url}
\usepackage{microtype}
\makeatletter\g@addto@macro\UrlBreaks{\do\_\do\-}\makeatother
\emergencystretch=2em

\newtheorem{theorem}{Theorem}[section]
\newtheorem{proposition}[theorem]{Proposition}
\newtheorem{lemma}[theorem]{Lemma}
\newtheorem{corollary}[theorem]{Corollary}
\newtheorem{conjecture}[theorem]{Conjecture}
\theoremstyle{definition}
\newtheorem{definition}[theorem]{Definition}
\newtheorem{remark}[theorem]{Remark}
\numberwithin{equation}{section}

\newcommand{\Z}{\mathbb{Z}}
\newcommand{\Q}{\mathbb{Q}}
\newcommand{\N}{\mathbb{N}}
\newcommand{\qbin}[2]{\genfrac{[}{]}{0pt}{}{#1}{#2}}
\newcommand{\repo}{\nolinkurl{2026-11-number-theory-q-series/}}
\newcommand{\fpath}[1]{\nolinkurl{#1}}
\newcommand{\T}{\mathsf T}

\title[Nahm sums for $E_6$, $E_7$, $E_8$]{Modular Nahm sums for the inverse Cartan matrices of
$E_6$, $E_7$ and $E_8$}
\author{Raghav Sharma}
\email{raghavsharma1729@gmail.com}
\thanks{\textbf{AI disclosure.} The results, proofs, computer code and Lean formalizations in this
paper were produced by AI (Anthropic's Claude models) working under the direction of the author.
Correctness was checked by independent AI referees (separate model instances instructed to find errors,
using their own code) and, for the formalized statements, by the Lean kernel. No human expert has
refereed this work yet. Code, certificates and Lean files:
\url{https://github.com/raghavs1729/Research-Math}, folder \texttt{2026-11-number-theory-q-series/}.}
\date{27 November 2026}

\subjclass[2020]{Primary 11P84; Secondary 05A30, 33D15, 11F03, 17B22, 81T40, 68V20}
\keywords{Nahm sums, Rogers--Ramanujan identities, Andrews--Gordon identities, inverse Cartan matrix,
exceptional root systems, Bailey's lemma, modular functions, Virasoro characters, Lean}

\begin{document}

\begin{abstract}
Cherednik and Feigin (2013) considered the Nahm sums
$\sum_{n}q^{n^{\T}An/2}/\prod_i(q)_{n_i}$ with $A$ the inverse Cartan matrix of a
simply-laced root system: they expected them to be modular (conditionally on Nahm's conjecture) and
noted that their method does not extend to this case. Sun and Wang (2026) observed
numerically that for $E_8$ and $E_6$ these sums are characters of the minimal models $M(2,11)$ and
$M(2,5)\otimes M(6,5)$, and found no formula for $E_7$. We prove both observations and give the $E_7$
formula: the Nahm sum of $C(E_7)^{-1}$ equals $\frac1{2(q;q)_\infty}\sum_{k\in\Z}(q^{3k^2/4}+q^{7k^2/4})$,
so it is modular. For $E_8$ we prove five identities modulo 11, including four companions identified by
Yoshida; for $E_6$ six identities, which rest on a theorem of Cao, Rosengren and Wang. The matrices
$C(E_8)^{-1}$, $C(E_6)^{-1}$, $C(E_7)^{-1}$ arise from small matrices by iterating a known Durfee-rectangle
fusion lemma, which reduces the sums to Andrews--Gordon, Rogers--Ramanujan and Zagier sums and, for
$E_7$, to a rank-four sum evaluated with Bailey's lemma and the $q$-Chu--Vandermonde sum. The same
argument gives an infinite family of identities. The matrix identities and fusion chains for $E_8$ and $E_6$ are checked in
Lean. An appendix records three misprints in Sun--Wang. The results, proofs, code and Lean files were
produced by AI (Anthropic's Claude models) directed by the author and checked by independent AI referees
and the Lean kernel; no human expert has refereed this work yet.
\end{abstract}

\maketitle

\section{Introduction}\label{sec:intro}

Throughout, $(a;q)_n=\prod_{j=0}^{n-1}(1-aq^j)$, $(q)_n=(q;q)_n$,
$(a_1,\dots,a_s;q)_\infty=\prod_i\prod_{j\ge0}(1-a_iq^j)$. For a symmetric matrix $A\in\Q^{r\times r}$ and
$B\in\Q^r$ put $Q_{A,B}(n)=\frac12n^{\T}An+B^{\T}n$ and
\begin{equation}\label{eq:F}
F_{A,B}(q)=\sum_{n\in\N^r}\frac{q^{Q_{A,B}(n)}}{(q)_{n_1}\cdots(q)_{n_r}} .
\end{equation}
Nahm's conjecture~\cite{Nahm,Zagier} relates the modularity of such sums (after multiplication by a
suitable $q^C$) to torsion in the Bloch group. For two Dynkin diagrams $X,Y$ of types $A,D,E,T$ the matrix
$C(X)\otimes C(Y)^{-1}$ is expected to give modular Nahm sums (a ``folklore'' conjecture, see Wang and
Wang~\cite[Conjecture~1.1]{WW-Dr}); Sun and Wang~\cite[Conjecture~1.1]{SW} extend it to all types with
symmetrizers and tabulate the conformal field theories that should explain them. For $X=T_1$, the tadpole
with $C(T_1)=(1)$, the matrix is $C(Y)^{-1}$, and Sun and Wang~\cite[§3]{SW} state, as relations found ``by
checking the $q$-series to very high orders'', identities equivalent to the following two (they write
the right sides as characters, see below):
\begin{align}
\sum_{n\in\N^8}\frac{q^{\frac12n^{\T}C(E_8)^{-1}n}}{\prod_i(q)_{n_i}}&=\frac1{(q,q^2,q^3,q^4,q^7,q^8,q^9,q^{10};q^{11})_\infty},
\tag{$\sim$SW 3.20}\label{eq:SW320}\\
\sum_{n\in\N^6}\frac{q^{\frac12n^{\T}C(E_6)^{-1}n}}{\prod_i(q)_{n_i}}&=\frac{J_{7,15}-qJ_{2,15}+2q^{2/3}J_{3,15}}{(q;q)_\infty\,(q,q^4;q^5)_\infty},
\tag{$\sim$SW 3.19}\label{eq:SW319}
\end{align}
with $J_{a,m}=(q^a,q^{m-a},q^m;q^m)_\infty$; i.e.\ the $(T_1,E_8)$ sum is the vacuum character of the
effective minimal model $M_{\rm eff}(11,2)$ (the Virasoro minimal model $M(2,11)$, with its characters
labelled by $h-h_{\min}$) and the $(T_1,E_6)$ sum is
$\chi_0^{M_{\rm eff}(5,2)}(\chi_0+\chi_3+2\chi_{2/3})^{M(6,5)}$. For $(T_1,E_7)$ they list the central charge
$c=1$ and the conformal field theory as ``?''~\cite[Table~1]{SW}. Much earlier, Cherednik and
Feigin~\cite[§3.6, ``Theta of level 1/2'']{CF} considered exactly these series: for a simply-laced root
system their ``flat'' matrix $A^\flat=((\omega_i,\omega_j))$ is $C(R)^{-1}$; they computed the
dilogarithm values numerically (their Table~1 gives $L^\flat=24/5,\,6,\,80/11$ for $E_6,E_7,E_8$, i.e.\
effective central charges $\mathrm{rank}-L^\flat=6/5,\,1,\,8/11$ in the usual normalisation of Nahm
sums; their own $c^\flat_{\rm eff}=L^\flat/h$ is normalised differently), regarded modularity in this case
as a conditional corollary of Nahm's conjecture, and wrote that their approach ``cannot
be directly extended to the $\flat$-case''.

We prove \eqref{eq:SW320} and \eqref{eq:SW319}, together with their companions, and find and prove the
missing $E_7$ formula.

\subsection{Main results}
Label $E_8$ as the chain $1-2-3-4-5-6-7$ with node $8$ attached to node $5$, $E_7$ as the chain
$1-2-3-4-5-6$ with node $7$ attached to node $4$, and $E_6$ as the chain $1-2-3-4-5$ with node $6$
attached to node $3$; $C(\cdot)$ is the Cartan matrix ($2$ on the diagonal, $-1$ for adjacent nodes).

\begin{theorem}[$E_8$]\label{thm:E8}
For every $b=(b_1,b_2,b_3,b_4)\in\Q^4$ let
$B(b)=(b_1,\,b_3,\,b_1+b_3,\,b_2+b_4,\,b_1+b_2+b_4,\,b_1+b_4,\,b_2,\,b_4)$. Then
\begin{equation}\label{eq:E8}
\sum_{n\in\N^8}\frac{q^{\frac12n^{\T}C(E_8)^{-1}n+B(b)\cdot n}}{\prod_{i=1}^8(q)_{n_i}}
=\sum_{m\in\N^4}\frac{q^{M_1^2+M_2^2+M_3^2+M_4^2+b\cdot m}}{(q)_{m_1}(q)_{m_2}(q)_{m_3}(q)_{m_4}},
\qquad M_j=m_j+\dots+m_4 .
\end{equation}
\end{theorem}

\begin{corollary}[five identities modulo 11]\label{cor:E8}
For $i=1,\dots,5$,
\[
\sum_{n\in\N^8}\frac{q^{\frac12n^{\T}C(E_8)^{-1}n+B^{(i)}\cdot n}}{\prod_{j=1}^8(q)_{n_j}}
=\prod_{\substack{m\ge1\\m\not\equiv0,\pm i\ (\mathrm{mod}\ 11)}}\frac1{1-q^m},
\]
with $B^{(5)}=0$, $B^{(4)}=(0,0,0,1,1,1,0,1)$, $B^{(3)}=(0,1,1,2,2,2,0,2)$, $B^{(2)}=(0,2,2,4,4,3,1,3)$,
$B^{(1)}=(1,3,4,6,7,5,2,4)$. The case $i=5$ is \eqref{eq:SW320}.
\end{corollary}

\begin{theorem}[$E_6$]\label{thm:E6}
For every $(b_1,b_2,b_3)\in\Q^3$ let $B=(b_1,\,b_2+b_3,\,b_1+b_2+b_3,\,b_1+b_3,\,b_2,\,b_3)$. Then
\[
\sum_{n\in\N^6}\frac{q^{\frac12n^{\T}C(E_6)^{-1}n+B\cdot n}}{\prod_{i=1}^6(q)_{n_i}}
=\Bigl(\sum_{i,j\ge0}\frac{q^{\frac23(i^2+ij+j^2)+b_1i+b_2j}}{(q)_i(q)_j}\Bigr)\Bigl(\sum_{m\ge0}\frac{q^{m^2+b_3m}}{(q)_m}\Bigr).
\]
\end{theorem}

\begin{corollary}[six identities for $E_6$]\label{cor:E6}
Let $G(q)=1/(q,q^4;q^5)_\infty$, $H(q)=1/(q^2,q^3;q^5)_\infty$, $J_m=(q^m;q^m)_\infty$,
\[
S_4(q)=\frac{J_{7,15}-qJ_{2,15}+2q^{2/3}J_{3,15}}{J_1},\qquad
S_3(q)=\frac{2J_{6,15}+q^{1/3}J_{4,15}+q^{4/3}J_{1,15}}{J_1}.
\]
Using the identities \eqref{eq:Ex10} for Zagier's rank-two Example~10, which are
Theorem~1.4 of Cao, Rosengren and Wang~\cite{CRW} (see §\ref{sec:E6}), $F_{C(E_6)^{-1},B}$ equals
\begin{center}
\begin{tabular}{ll}
$G(q)S_4(q)$ for $B=(0,0,0,0,0,0)$, & $H(q)S_4(q)$ for $B=(0,1,1,1,0,1)$,\\
$G(q)S_3(q)$ for $B=(-\frac23,-\frac13,-1,-\frac23,-\frac13,0)$, & $H(q)S_3(q)$ for $B=(-\frac23,\frac23,0,\frac13,-\frac13,1)$,\\
$G(q)S_3(q)$ for $B=(-\frac13,-\frac23,-1,-\frac13,-\frac23,0)$, & $H(q)S_3(q)$ for $B=(-\frac13,\frac13,0,\frac23,-\frac23,1)$.
\end{tabular}
\end{center}
The first is \eqref{eq:SW319}.
\end{corollary}

\begin{theorem}[$E_7$]\label{thm:E7}
\[
\sum_{n\in\N^7}\frac{q^{\frac12n^{\T}C(E_7)^{-1}n}}{\prod_{i=1}^7(q)_{n_i}}
=\frac1{(q;q)_\infty}\Bigl(\sum_{m\ge0}q^{3m^2/4}+\sum_{m\ge1}q^{7m^2/4}\Bigr)
=\frac1{2(q;q)_\infty}\sum_{m\in\Z}\bigl(q^{3m^2/4}+q^{7m^2/4}\bigr).
\]
Consequently $q^{-1/24}$ times the left side is $(\theta(3\tau/4)+\theta(7\tau/4))/(2\eta(\tau))$, where
$\theta(z)=\sum_{m\in\Z}e^{2\pi im^2z}$ and $q=e^{2\pi i\tau}$, a modular function for a congruence
subgroup; so Conjecture~1.1 of~\cite{SW} holds for the pair $(T_1,E_7)$.
\end{theorem}

Separating the exponents in $\Z$ and in $\frac34+\Z$ gives two refinements (§\ref{sec:E7}): the parts
of the sum with $n_1+n_3+n_7$ even, resp.\ odd, equal
\[
\frac1{2(q)_\infty}\sum_{k\in\Z}\bigl(q^{3k^2}+q^{7k^2}\bigr),\qquad\text{resp.}\qquad
\frac1{2(q)_\infty}\sum_{k\in\Z}\bigl(q^{3(k+1/2)^2}+q^{7(k+1/2)^2}\bigr).
\] The right side of
Theorem~\ref{thm:E7} is the average of the $(A_1,D_3)$ and $(A_1,D_7)$ Nahm sums
$\sum_nq^{n^{\T}C(D_r)^{-1}n}/\prod(q)_{n_i}=(q)_\infty^{-1}\sum_kq^{rk^2/4}$
(Kedem, Klassen, McCoy and Melzer~\cite{KKMM1}; see~\cite[Theorem~1.2]{WW-Dk} and the references
there).

The proof of Theorem~\ref{thm:E7} gives an infinite family.

\begin{theorem}\label{thm:family}
For every integer $k\ge1$, with $N_j=n_j+\dots+n_k$,
\[
\sum_{a,n_1,\dots,n_k\ge0}\frac{q^{\frac k4a^2+a(N_2+\dots+N_k)+N_1^2+\dots+N_k^2}}{(q)_a(q)_{n_1}\cdots(q)_{n_k}}
=\frac1{(q)_\infty}\Bigl(\sum_{m\ge0}q^{km^2/4}+\sum_{m\ge1}q^{(k+4)m^2/4}\Bigr).
\]
\end{theorem}

\subsection{What is new, what was known, what is open}\label{ssec:novelty}
\emph{Known.} The fusion lemma (Lemma~\ref{lem:fusion}), in arbitrary rank and for an arbitrary pair of
indices, is in print: Cao and Wang~\cite[(1.5)--(1.6)]{CW3} (``lifting''; also~\cite{CW-lift3,CW-lift4}),
and, as the ``linking'' of symmetric quivers, Ekholm, Kucharski and Longhi~\cite[§4.5,
(4.31)--(4.36)]{EKL}, iterated systematically by Kucharski, Larragu\'ivel, Noshchenko and
Su\l kowski~\cite{KLNS}; it is also used by L.~Wang~\cite{Wang22} and by Cao, Rosengren and
Wang~\cite[(2.4)]{CRW}. Lifts reaching an inverse Cartan matrix occur for $C(D_3)^{-1}=C(A_3)^{-1}$
(\cite[§4]{WW-Dk}, \cite[§2]{SW}). The Durfee identity, Bailey's lemma and the $q$-Chu--Vandermonde sum
are classical. The case $k=2$ of Theorem~\ref{thm:family} is the $(T_1,T_4)$ identity, the case of four
summation variables and no linear term of Melzer's conjecture~\cite[(2.6)]{Melzer}, proved by Bressoud,
Ismail and Stanton~\cite[Theorem~5.1]{BIS} and by Warnaar~\cite[Theorem~4.1]{Warnaar03} (Theorem~4.4 in
the numbering quoted in~\cite{SW}); our proof of it is
new but the identity is not. The rank-two identities used for $E_6$ are theorems of Cao, Rosengren and
Wang~\cite{CRW}. The case $k=1$ of Theorem~\ref{thm:family} also follows from known identities: the left
side is $G(q)\sum_aq^{a^2/4}/(q)_a$, and Rogers' identity
$\sum_aq^{a^2/4}/(q)_a=G(q^{1/4})/(-q^{1/2};q^{1/2})_\infty$ together with Ramanujan's
$G(p)G(p^4)+pH(p)H(p^4)=\varphi(p)/(p^2;p^2)_\infty$ and $G(p)G(p^4)-pH(p)H(p^4)=\varphi(p^5)/(p^2;p^2)_\infty$
($p=q^{1/4}$, $\varphi(p)=\sum_{m\in\Z}p^{m^2}$; two of Ramanujan's forty identities, see~\cite{Forty})
gives the right side. Yoshida~\cite{Yoshida} identified the four companions $i\le4$ of
Corollary~\ref{cor:E8} (his charges are $-B^{(i)}$) and verified them ``through relative order
$q^6$''; he treats the vacuum identity \eqref{eq:SW320} as known, citing only~\cite{SW}.

\emph{New.} The observation that $C(E_8)^{-1}$, $C(E_6)^{-1}$ and $C(E_7)^{-1}$ are iterated lifts of
$2C(T_4)^{-1}$, of $2C(A_2)^{-1}\oplus(2)$ and of a rank-four matrix, and its consequences:
Theorems~\ref{thm:E8} and~\ref{thm:E6}, Corollary~\ref{cor:E8} (hence \eqref{eq:SW320} and Yoshida's
companions), Corollary~\ref{cor:E6} given~\cite{CRW} (hence \eqref{eq:SW319}), the formula of
Theorem~\ref{thm:E7} and the modularity of the $(T_1,E_7)$ sum, and Theorem~\ref{thm:family} for
$k\ge3$ (and the new proofs for $k=1,2$). As far as our searches go, no proof of \eqref{eq:SW320} or \eqref{eq:SW319} and no formula for
$(T_1,E_7)$ existed.

\emph{Open.} The dual pairs $(E_n,T_1)$, i.e.\ the Nahm sums of $C(E_n)$ itself (see
§\ref{sec:open}); the other unidentified rows of~\cite[Table~1]{SW}.

\emph{Searches} (26 November 2026): arXiv metadata searches up to 25 November 2026 (``Nahm sum(s)'' with
$E_6$, $E_7$, $E_8$, Dynkin, inverse Cartan; Rogers--Ramanujan and Andrews--Gordon with $E_8$; recent
work of the authors in this area); the papers citing~\cite{SW} (Semantic Scholar, INSPIRE) and
the papers citing~\cite{CF}; full texts of the most relevant papers, among
them~\cite{SW,CF,CRW,CW3,WW-Dr,WW-Dk,Yoshida,KKMM1,KKMM2,WP,VZ,Nahm,Melzer,BIS,Warnaar03}. The
1990s fermionic-character literature concerns $2C(G)^{-1}$ (for $E_8$: the Ising model) rather than
$C(G)^{-1}$~\cite{KKMM1,KKMM2,WP}. No full-text search engine (MathSciNet, zbMATH, Google Scholar) was
available; since the proofs are short and use known tools, an unindexed independent proof cannot be
excluded. An anticipation of Theorem~\ref{thm:family} in the Bailey-lattice literature was not checked.

\subsection{Formalization}
The finite and algebraic ingredients are checked in Lean~4 with Mathlib~\cite{Lean4,Mathlib}: the printed
matrices are $C(E_8)^{-1}$ and $C(E_6)^{-1}$; the fusion chains produce exactly these matrices and the
stated linear terms; the quadratic-form identity behind the fusion lemma; and the Durfee identity
(§\ref{sec:lean}). The infinite-series identities themselves are not yet formalized; a separate effort
towards this is in progress.

\subsection{Organisation}
§\ref{sec:lemmas}: setting and the two lemmas. §§\ref{sec:E8}, \ref{sec:E6}, \ref{sec:E7}: the three
root systems. §\ref{sec:family}: Theorem~\ref{thm:family}. §\ref{sec:open}: further remarks and open
problems. §\ref{sec:lean}: Lean. §\ref{sec:num}: numerical checks. Appendix~\ref{app:SW}: remarks on~\cite{SW}.

\section{Setting and two lemmas}\label{sec:lemmas}

Call $(A,B)$ \emph{admissible} if $\{n\in\N^r:Q_{A,B}(n)\le E\}$ is finite for every real $E$. Then, with
$D$ a common denominator of the entries of $A/2$ and of $B$, \eqref{eq:F} is a summable family in the ring
$\Z((q^{1/D}))$ of formal Laurent series (each term is $q^{Q(n)}(1+O(q))$ with integer coefficients), and
$F_{A,B}$ is well defined there. If $A$ is positive definite, $(A,B)$ is admissible for every $B$. All
identities below are identities in $\Z((q^{1/D}))$; both sides also converge for $|q|<1$.

\begin{lemma}[Durfee rectangle identity]\label{lem:durfee}
For integers $a,b\ge0$,
\begin{equation}\label{eq:durfee}
\frac1{(q)_a(q)_b}=\sum_{c=0}^{\min(a,b)}\frac{q^{(a-c)(b-c)}}{(q)_{a-c}(q)_{b-c}(q)_c}.
\end{equation}
\end{lemma}
\begin{proof}
Work in $\Z[[q]][[x,y]]$. By Euler, $\sum_cz^c/(q)_c=1/(z;q)_\infty$, and by the $q$-binomial theorem
$\sum_i(y;q)_ix^i/(q)_i=(xy;q)_\infty/(x;q)_\infty$ \cite[Theorem~2.1, (2.2.5)]{AndrewsBook} (both
sides are power series in the variables whose coefficients are rational functions of $q$, resp.\ $q$ and
$y$, agreeing on an open set, hence identities of formal power series). Hence
\[
\sum_{i,j,c\ge0}\frac{q^{ij}x^{i+c}y^{j+c}}{(q)_i(q)_j(q)_c}
=\frac1{(xy;q)_\infty}\sum_i\frac{x^i}{(q)_i}\frac1{(yq^i;q)_\infty}
=\frac1{(xy;q)_\infty(y;q)_\infty}\sum_i\frac{(y;q)_ix^i}{(q)_i}=\frac1{(x;q)_\infty(y;q)_\infty},
\]
and comparing the coefficients of $x^ay^b$ ($a=i+c$, $b=j+c$) gives \eqref{eq:durfee}.
\end{proof}

\begin{lemma}[fusion]\label{lem:fusion}
Let $(A,B)$ be admissible of rank $r$ and let $a\ne b$ be two indices. Define rank $r+1$ data
$(A^\circ,B^\circ)$, with an additional index $0$, by
\begin{gather*}
A^\circ_{ij}=A_{ij}+[\{i,j\}=\{a,b\}]\ (1\le i,j\le r),\qquad A^\circ_{i0}=A^\circ_{0i}=A_{ia}+A_{ib}\ (1\le i\le r),\\
A^\circ_{00}=A_{aa}+A_{bb}+2A_{ab},\qquad B^\circ_i=B_i\ (1\le i\le r),\qquad B^\circ_0=B_a+B_b .
\end{gather*}
Here $[P]=1$ if $P$ holds and $0$ otherwise. Then $(A^\circ,B^\circ)$ is admissible and
$F_{A^\circ,B^\circ}=F_{A,B}$.
\end{lemma}
\begin{proof}
Let $v=e_a+e_b$. For $(m,c)\in\N^r\times\N$,
\begin{equation}\label{eq:Qfuse}
Q_{A^\circ,B^\circ}(m,c)=\tfrac12m^{\T}Am+m_am_b+c\,(Av)^{\T}m+\tfrac12c^2v^{\T}Av+B^{\T}m+cB^{\T}v=Q_{A,B}(m+cv)+m_am_b .
\end{equation}
The map $(m,c)\mapsto n=m+cv$ has finite fibres ($c\le\min(n_a,n_b)$), and $Q_{A^\circ,B^\circ}(m,c)\ge
Q_{A,B}(n)$, so $(A^\circ,B^\circ)$ is admissible. Grouping the summable family $F_{A^\circ,B^\circ}$
according to $n$ and using \eqref{eq:Qfuse},
\[
F_{A^\circ,B^\circ}=\sum_{n\in\N^r}\frac{q^{Q_{A,B}(n)}}{\prod_{i\ne a,b}(q)_{n_i}}
\sum_{c=0}^{\min(n_a,n_b)}\frac{q^{(n_a-c)(n_b-c)}}{(q)_{n_a-c}(q)_{n_b-c}(q)_c}=F_{A,B}
\]
by Lemma~\ref{lem:durfee}.
\end{proof}

In words: one may adjoin a species $0=a+b$ whose row of $A$ is the sum of the rows of $a$ and $b$, at the
cost of raising $A_{ab}$ by $1$. In matrix form $A^\circ=P^{\T}AP+E_{ab}+E_{ba}$, $B^\circ=P^{\T}B$ with
$P=(I_r\,|\,v)$. Positive definiteness is not needed, although all intermediate matrices below are in fact
positive definite. As explained in §\ref{ssec:novelty}, Lemma~\ref{lem:fusion} is known; it is
included to keep the paper self-contained.

\emph{Formal-variable version.} If $A$ is positive definite and $y_1,\dots,y_r$ are indeterminates,
then
$\sum_nq^{\frac12n^{\T}An}y^n/\prod(q)_{n_i}=\sum_{(m,c)}q^{\frac12m^{\T}Am+m_am_b+c(Av)^{\T}m+\frac12c^2v^{\T}Av}y^m(y_ay_b)^c/(\prod(q)_{m_i}\cdot(q)_c)$
in $\Z[[q^{1/D}]][[y_1,\dots,y_r]]$: the coefficient of each monomial $y^n$ is a finite sum ($c\le\min(n_a,n_b)$),
equal on both sides by Lemma~\ref{lem:durfee}.

\section{\texorpdfstring{$E_8$}{E8}}\label{sec:E8}

\begin{proof}[Proof of Theorem~\ref{thm:E8}]
Let $A^{(0)}=(2\min(i,j))_{1\le i,j\le4}$; then $\frac12m^{\T}A^{(0)}m=\sum_{i,j}\min(i,j)m_im_j=M_1^2+\dots+M_4^2$,
so the right side of \eqref{eq:E8} is $F_{A^{(0)},b}$ ($A^{(0)}=2C(T_4)^{-1}$ is positive definite). Name the
four variables $m_1,m_2,m_3,m_4$ after the $E_8$ nodes $1,7,2,8$, and apply Lemma~\ref{lem:fusion} four
times, creating the nodes $4,6,5,3$ in this order:
\[
\text{node }4:=7+8,\qquad\text{node }6:=1+8,\qquad\text{node }5:=1+4,\qquad\text{node }3:=1+2 .
\]
Each step copies ``row $a$ + row $b$'' into the new row, puts $A_{aa}+A_{bb}+2A_{ab}$ on the diagonal, and
then raises $A_{ab}$ by $1$. With rows and columns in the order of creation, the five stages are
\[
\begin{array}{c|cccc}
&1&7&2&8\\\hline
1&2&2&2&2\\7&2&4&4&4\\2&2&4&6&6\\8&2&4&6&8
\end{array}
\qquad
\begin{array}{c|ccccc}
&1&7&2&8&4\\\hline
1&2&2&2&2&4\\7&2&4&4&5&8\\2&2&4&6&6&10\\8&2&5&6&8&12\\4&4&8&10&12&20
\end{array}
\qquad
\begin{array}{c|cccccc}
&1&7&2&8&4&6\\\hline
1&2&2&2&3&4&4\\7&2&4&4&5&8&7\\2&2&4&6&6&10&8\\8&3&5&6&8&12&10\\4&4&8&10&12&20&16\\6&4&7&8&10&16&14
\end{array}
\]
\[
\begin{array}{c|ccccccc}
&1&7&2&8&4&6&5\\\hline
1&2&2&2&3&5&4&6\\7&2&4&4&5&8&7&10\\2&2&4&6&6&10&8&12\\8&3&5&6&8&12&10&15\\
4&5&8&10&12&20&16&24\\6&4&7&8&10&16&14&20\\5&6&10&12&15&24&20&30
\end{array}
\qquad
\begin{array}{c|cccccccc}
&1&7&2&8&4&6&5&3\\\hline
1&2&2&3&3&5&4&6&4\\7&2&4&4&5&8&7&10&6\\2&3&4&6&6&10&8&12&8\\8&3&5&6&8&12&10&15&9\\
4&5&8&10&12&20&16&24&15\\6&4&7&8&10&16&14&20&12\\5&6&10&12&15&24&20&30&18\\3&4&6&8&9&15&12&18&12
\end{array}
\]
(the raised entries are $A_{78}:4\to5$, $A_{18}:2\to3$, $A_{14}:4\to5$, $A_{12}:2\to3$). Reordered as
$1,\dots,8$, the last matrix is
\[
\begin{pmatrix}
2&3&4&5&6&4&2&3\\3&6&8&10&12&8&4&6\\4&8&12&15&18&12&6&9\\5&10&15&20&24&16&8&12\\
6&12&18&24&30&20&10&15\\4&8&12&16&20&14&7&10\\2&4&6&8&10&7&4&5\\3&6&9&12&15&10&5&8
\end{pmatrix}=C(E_8)^{-1},
\]
as one checks by multiplying by $C(E_8)$. Under the four steps the linear term $b$ on the nodes
$(1,7,2,8)$ acquires the entries $B_4=b_2+b_4$, $B_6=b_1+b_4$, $B_5=b_1+b_2+b_4$, $B_3=b_1+b_3$, which is the
vector $B(b)$. By Lemma~\ref{lem:fusion}, $F_{C(E_8)^{-1},B(b)}=F_{A^{(0)},b}$.
\end{proof}

\begin{proof}[Proof of Corollary~\ref{cor:E8}]
The Andrews--Gordon identities~\cite[Theorem~1]{Andrews74} (\cite[Theorem~7.8]{AndrewsBook};
Gordon~\cite{Gordon} for the combinatorial version) state that for $k\ge2$ and $1\le i\le k$,
\[
\sum_{m_1,\dots,m_{k-1}\ge0}\frac{q^{M_1^2+\dots+M_{k-1}^2+M_i+\dots+M_{k-1}}}{(q)_{m_1}\cdots(q)_{m_{k-1}}}
=\prod_{\substack{n\ge1\\n\not\equiv0,\pm i\ (\mathrm{mod}\ 2k+1)}}\frac1{1-q^n},\qquad M_j=m_j+\dots+m_{k-1}.
\]
Take $k=5$. Since $M_i+\dots+M_4=\sum_j\max(0,j-i+1)m_j$, the left side is $F_{A^{(0)},b^{(i)}}$ with
$b^{(5)}=0$, $b^{(4)}=(0,0,0,1)$, $b^{(3)}=(0,0,1,2)$, $b^{(2)}=(0,1,2,3)$, $b^{(1)}=(1,2,3,4)$, and
Theorem~\ref{thm:E8} gives the corollary with $B^{(i)}=B(b^{(i)})$.
\end{proof}

\emph{Modularity and the CFT reading.} The characters of the Virasoro minimal model $M(2,11)$ are
$q^{h_i-c/24}\prod_{n\not\equiv0,\pm i\ (11)}(1-q^n)^{-1}$, $i=1,\dots,5$, with $c=-232/11$ and
$h_i=-(i-1)(10-i)/22$ (Rocha-Caridi's formula~\cite{RC} and the Jacobi triple product); the lowest weight
is $h_5$, and $h_5-c/24=-1/33$. With $C=-c_{\rm eff}/24=-1/33$, $c_{\rm eff}=8/11$, the case $i=5$ reads
$q^{-1/33}F_{C(E_8)^{-1},0}=\eta(\tau)^{-1}\sum_{n\in\Z}(-1)^nq^{11(n+1/22)^2/2}$ (Jacobi triple product;
$-\frac1{33}-\frac1{88}+\frac1{24}=0$), a weight-$\frac12$ unary theta function divided by $\eta$, hence a
modular function for a congruence subgroup; this establishes Conjecture~1.1 of~\cite{SW} (and the
folklore conjecture~\cite[Conjecture~1.1]{WW-Dr}) for $(T_1,E_8)$. Of the five products only $i=5$
agrees with the sum for $B=0$ (the products $i=1,2,3,4$ first differ from it at $q^1,q^2,q^3,q^4$). The
effective weights $h_i-h_5$ are $10/11,6/11,3/11,1/11$ for $i=1,2,3,4$, in agreement with Yoshida's
identifications~\cite[(3.77)--(3.80)]{Yoshida}.

\begin{remark}[refined form]
The same computation with indeterminates gives, in $\Z[[q]][[y_1,\dots,y_4]]$,
\[
\sum_{n\in\N^8}\frac{q^{\frac12n^{\T}C(E_8)^{-1}n}y_1^{n_1+n_3+n_5+n_6}y_2^{n_4+n_5+n_7}y_3^{n_2+n_3}y_4^{n_4+n_5+n_6+n_8}}{\prod(q)_{n_i}}
=\sum_{m\in\N^4}\frac{q^{M_1^2+\dots+M_4^2}y_1^{m_1}y_2^{m_2}y_3^{m_3}y_4^{m_4}}{\prod(q)_{m_j}} .
\]
So the eight $E_8$ ``quasi-particles'' are the four Andrews--Gordon ones (nodes $1,7,2,8$) and the bound
states $1+2$, $1+8$, $7+8$, $1+7+8$ (nodes $3,6,4,5$). With $y_j=x^j$ the weight is $x^{w\cdot n}$,
$w=C(E_8)^{-1}\varepsilon=(1,3,4,6,7,5,2,4)$, $\varepsilon=(-1,1,-1,1,-1,1,-1,1)$ the bipartition sign vector,
and the right side is the function $J_{5,5}(0;x;q)$ of~\cite[Ch.~7]{AndrewsBook}.
\end{remark}

\section{\texorpdfstring{$E_6$}{E6}}\label{sec:E6}

\begin{proof}[Proof of Theorem~\ref{thm:E6}]
Start with the rank-three data
$A^{(0)}=\left(\begin{smallmatrix}4/3&2/3&0\\2/3&4/3&0\\0&0&2\end{smallmatrix}\right)=2C(A_2)^{-1}\oplus(2)$, $b=(b_1,b_2,b_3)$,
whose Nahm sum is the right side of the theorem (it factors because $A^{(0)}$ is block diagonal). Name
the variables $(i,j,m)$ after the $E_6$ nodes $1,5,6$ and apply Lemma~\ref{lem:fusion} three times:
node $2:=5+6$, node $4:=1+6$, node $3:=1+2$. In the order $(1,5,6\,|\,2\,|\,4\,|\,3)$ the stages are
\[
\begin{pmatrix}\frac43&\frac23&0\\\frac23&\frac43&0\\0&0&2\end{pmatrix}\!,
\begin{pmatrix}\frac43&\frac23&0&\frac23\\\frac23&\frac43&1&\frac43\\0&1&2&2\\\frac23&\frac43&2&\frac{10}3\end{pmatrix}\!,
\begin{pmatrix}\frac43&\frac23&1&\frac23&\frac43\\\frac23&\frac43&1&\frac43&\frac53\\1&1&2&2&2\\\frac23&\frac43&2&\frac{10}3&\frac83\\\frac43&\frac53&2&\frac83&\frac{10}3\end{pmatrix}\!,
\begin{pmatrix}\frac43&\frac23&1&\frac53&\frac43&2\\\frac23&\frac43&1&\frac43&\frac53&2\\1&1&2&2&2&3\\\frac53&\frac43&2&\frac{10}3&\frac83&4\\\frac43&\frac53&2&\frac83&\frac{10}3&4\\2&2&3&4&4&6\end{pmatrix}\!,
\]
and the last matrix, reordered as $1,\dots,6$, is
\[
\begin{pmatrix}
4/3&5/3&2&4/3&2/3&1\\5/3&10/3&4&8/3&4/3&2\\2&4&6&4&2&3\\4/3&8/3&4&10/3&5/3&2\\2/3&4/3&2&5/3&4/3&1\\1&2&3&2&1&2
\end{pmatrix}=C(E_6)^{-1}.
\]
The linear terms become $B_2=b_2+b_3$, $B_4=b_1+b_3$, $B_3=b_1+b_2+b_3$. Lemma~\ref{lem:fusion} gives the
theorem.
\end{proof}

For Corollary~\ref{cor:E6} we need the Rogers--Ramanujan identities
$\sum_mq^{m^2}/(q)_m=G(q)$, $\sum_mq^{m^2+m}/(q)_m=H(q)$~\cite{Rogers,AndrewsBook} and
\begin{equation}\label{eq:Ex10}
\sum_{i,j\ge0}\frac{q^{\frac23(i^2+ij+j^2)}}{(q)_i(q)_j}=S_4(q),\qquad
\sum_{i,j\ge0}\frac{q^{\frac23(i^2+ij+j^2)-\frac{2i+j}3}}{(q)_i(q)_j}=S_3(q)
\end{equation}
(and the same with $i,j$ exchanged). These are the identities for Zagier's rank-two Example~10,
$A=\left(\begin{smallmatrix}4/3&2/3\\2/3&4/3\end{smallmatrix}\right)$, conjectured by Vlasenko and
Zwegers~\cite[Table~1]{VZ}; under $q\mapsto q^3$ they are equations (1.6b) and (1.6a) of Cao, Rosengren and
Wang~\cite{CRW}. Their status, precisely: Cherednik and Feigin~\cite[around (2.20)]{CF} prove modularity of
the sum side for $B=0$ and report a computer comparison of coefficients (no bound stated); L.~Wang~\cite[§3.10]{Wang22}
still lists both as conjectures; Cao, Rosengren and Wang prove $S_1=(-q^3;q^3)_\infty S_3$,
$S_2=(-q^3;q^3)_\infty S_4$ (in their normalisation) and product formulas for $S_1,S_2$, and deduce
\eqref{eq:Ex10} from these by an identity between modular theta quotients verified with a Maple
package~\cite[Theorem~1.4]{CRW}. \emph{Corollary~\ref{cor:E6} and \eqref{eq:SW319} are therefore exactly as
rigorous as~\cite[Theorem~1.4]{CRW}; Theorem~\ref{thm:E6} itself is unconditional.}

\begin{proof}[Proof of Corollary~\ref{cor:E6}]
Insert the Rogers--Ramanujan identities and \eqref{eq:Ex10} in Theorem~\ref{thm:E6} with
$(b_1,b_2)\in\{(0,0),(-\frac23,-\frac13),(-\frac13,-\frac23)\}$ and $b_3\in\{0,1\}$.
\end{proof}

\emph{The CFT form of \eqref{eq:SW319}.} We have $\chi_0^{M_{\rm eff}(5,2)}=q^{-1/60}G(q)$, and
$q^{-1/30}S_4=(\chi_0+\chi_3+2\chi_{2/3})^{M(6,5)}$. For the latter, by the Jacobi triple product
$J_{a,15}=\sum_n(-1)^nq^{(15n^2-(15-2a)n)/2}$; splitting $n=2k,2k+1$ gives
\begin{align*}
J_{7,15}&=\textstyle\sum_k(q^{30k^2-k}-q^{30k^2+29k+7}),\qquad
-qJ_{2,15}=\sum_k(-q^{30k^2-11k+1}+q^{30k^2+19k+3}),\\
J_{3,15}&=\textstyle\sum_k(q^{30k^2-9k}-q^{30k^2+21k+3}).
\end{align*}
The Rocha-Caridi numerators of $M(6,5)$ (with $p=5$, $p'=6$) are $N_{r,s}=\sum_k(q^{k(30k+6r-5s)}-q^{(5k+r)(6k+s)})$, so
$N_{1,1}=\sum_k(q^{30k^2+k}-q^{30k^2+11k+1})$, $q^3N_{4,1}=\sum_k(q^{30k^2+19k+3}-q^{30k^2+29k+7})$ and
$N_{1,3}=\sum_k(q^{30k^2-9k}-q^{30k^2+21k+3})$. Hence $J_{3,15}=N_{1,3}$, and
$J_{7,15}-qJ_{2,15}=N_{1,1}+q^3N_{4,1}$: the four half-series of the left side coincide with the four
half-series of the right side, paired crosswise (with $k\mapsto-k$ in $q^{30k^2-k}$ and $q^{30k^2-11k+1}$);
individually $J_{7,15}\ne N_{1,1}$. With $h=0,3,2/3$ for $(r,s)=(1,1),(4,1),(1,3)$ this is the claimed
identification. Each of $q^{-1/60}G$ and $q^{-1/30}S_4$ is a combination of weight-$\frac12$ unary theta
functions divided by $\eta$, so $q^{-1/20}F_{C(E_6)^{-1},0}$ (with $-\frac1{20}=-\frac1{60}-\frac1{30}=-c_{\rm eff}/24$,
$c_{\rm eff}=6/5$) is modular and Conjecture~1.1 of~\cite{SW} holds for $(T_1,E_6)$,
given~\cite[Theorem~1.4]{CRW}.

\section{\texorpdfstring{$E_7$}{E7}}\label{sec:E7}

\subsection{Reduction to rank four}
Let $A^{(0)}$ be the matrix on the variables $(a,n_1,n_2,n_3)$, named after the $E_7$ nodes $1,6,2,5$:
\[
A^{(0)}=\begin{pmatrix}3/2&0&1&2\\0&2&2&2\\1&2&4&4\\2&2&4&6\end{pmatrix},\qquad
\tfrac12(a,n)A^{(0)}(a,n)^{\T}=\tfrac34a^2+a(N_2+N_3)+N_1^2+N_2^2+N_3^2,
\]
$N_j=n_j+\dots+n_3$ ($A^{(0)}$ is positive definite, $\det A^{(0)}=4$). Apply Lemma~\ref{lem:fusion} three
times: node $7:=1+6$, node $4:=5+6$, node $3:=1+2$. In the order $(1,6,2,5\,|\,7\,|\,4\,|\,3)$ the stages are
\[
\begin{pmatrix}\frac32&1&1&2&\frac32\\1&2&2&2&2\\1&2&4&4&3\\2&2&4&6&4\\\frac32&2&3&4&\frac72\end{pmatrix}\!,\
\begin{pmatrix}\frac32&1&1&2&\frac32&3\\1&2&2&3&2&4\\1&2&4&4&3&6\\2&3&4&6&4&8\\\frac32&2&3&4&\frac72&6\\3&4&6&8&6&12\end{pmatrix}\!,\
\begin{pmatrix}\frac32&1&2&2&\frac32&3&\frac52\\1&2&2&3&2&4&3\\2&2&4&4&3&6&5\\2&3&4&6&4&8&6\\\frac32&2&3&4&\frac72&6&\frac92\\3&4&6&8&6&12&9\\\frac52&3&5&6&\frac92&9&\frac{15}2\end{pmatrix}
\]
(raised entries $A_{16}:0\to1$, $A_{56}:2\to3$, $A_{12}:1\to2$), and the last matrix reordered as
$1,\dots,7$ is
\[
\begin{pmatrix}
3/2&2&5/2&3&2&1&3/2\\2&4&5&6&4&2&3\\5/2&5&15/2&9&6&3&9/2\\3&6&9&12&8&4&6\\2&4&6&8&6&3&4\\1&2&3&4&3&2&2\\3/2&3&9/2&6&4&2&7/2
\end{pmatrix}=C(E_7)^{-1}.
\]
Hence, by Lemma~\ref{lem:fusion},
\begin{equation}\label{eq:E7red}
F(q):=\sum_{n\in\N^7}\frac{q^{\frac12n^{\T}C(E_7)^{-1}n}}{\prod(q)_{n_i}}
=\sum_{a\ge0}\frac{q^{3a^2/4}}{(q)_a}\,T_3(a),\qquad
T_k(a):=\sum_{n_1,\dots,n_k\ge0}\frac{q^{N_1^2+\dots+N_k^2+a(N_2+\dots+N_k)}}{(q)_{n_1}\cdots(q)_{n_k}} .
\end{equation}
Thus Theorem~\ref{thm:E7} is the case $k=3$ of Theorem~\ref{thm:family}, which we now prove.

\section{Proof of Theorem~\ref{thm:family}}\label{sec:family}

\begin{lemma}[Bailey chain]\label{lem:bailey}
Let $x$ be a parameter with $x\ne q^{-m}$ $(m\ge1)$ and put, for $r\ge0$,
\[
\alpha_r=\frac{(-1)^rq^{r(r-1)/2}(1-xq^{2r})(x;q)_r}{(1-x)(q)_r}
\]
($\alpha_0=1$; for $x=1$ read $(x;q)_r/(1-x)$ as $(q;q)_{r-1}$, so $\alpha_r=(-1)^rq^{r(r-1)/2}(1+q^r)$ for $r\ge1$).
Then for every $k\ge1$ and $N\ge0$
\begin{equation}\label{eq:bailey}
\sum_{N\ge N_2\ge\dots\ge N_k\ge0}\frac{x^{N_2+\dots+N_k}q^{N_2^2+\dots+N_k^2}}{(q)_{N-N_2}(q)_{N_2-N_3}\cdots(q)_{N_{k-1}-N_k}(q)_{N_k}}
=\sum_{r=0}^{N}\frac{x^{kr}q^{kr^2}\alpha_r}{(q)_{N-r}(xq;q)_{N+r}},
\end{equation}
where for $k=1$ the left side is $1/(q)_N$.
\end{lemma}
\begin{proof}
A pair $(\alpha_n,\beta_n)$ is a Bailey pair relative to $x$ if
$\beta_n=\sum_{r=0}^n\alpha_r/((q)_{n-r}(xq)_{n+r})$. (i) The unit Bailey pair is $\alpha_r$ as above,
$\beta_n=\delta_{n,0}$. (ii) Bailey's lemma in the limit $\rho_1,\rho_2\to\infty$: if $(\alpha,\beta)$ is a
Bailey pair relative to $x$, then so is $\alpha'_r=x^rq^{r^2}\alpha_r$,
$\beta'_n=\sum_{r=0}^nx^rq^{r^2}\beta_r/(q)_{n-r}$. Both are classical: Andrews~\cite[(2.3)--(2.6), Theorem~1
(2.7), and (2.12)--(2.13) for the unit pair]{Andrews84}; Warnaar~\cite[(1.3), Theorem~2.1, (2.2), (2.3)]{Warnaar50}.
Apply (ii) $k$ times to (i). By induction on $k$, $\alpha^{(k)}_r=x^{kr}q^{kr^2}\alpha_r$;
$\beta^{(1)}_n=1/(q)_n$; and if $\beta^{(k)}_n$ is the left side of \eqref{eq:bailey} with $N=n$, then
$\beta^{(k+1)}_N=\sum_{N_2\le N}x^{N_2}q^{N_2^2}\beta^{(k)}_{N_2}/(q)_{N-N_2}$ is the left side for $k+1$ (rename the
summation variables of $\beta^{(k)}$ as $N_3,\dots,N_{k+1}$). Identity \eqref{eq:bailey} is the statement that
$(\alpha^{(k)},\beta^{(k)})$ is a Bailey pair. Both sides are rational functions of $x$ and $q$, regular at
$x=q^a$, $a\ge0$.
\end{proof}

\begin{lemma}\label{lem:cauchy}
For integers $a,R\ge0$,
\[
\sum_{m\ge0}\frac{q^{m^2+Rm}}{(q)_m(q)_{m+R+a}}=\frac1{(q)_\infty}\sum_{j=0}^{a}\qbin aj q^{j^2+Rj},
\qquad\qbin aj=\frac{(q)_a}{(q)_j(q)_{a-j}} .
\]
\end{lemma}
\begin{proof}
Cauchy's identity $\sum_{m\ge0}q^{m(m+c)}/((q)_m(q)_{m+c})=1/(q)_\infty$ holds for every integer $c\ge0$: Lemma~\ref{lem:durfee}
with $(a,b)\to(A,A+c)$ reads $1/((q)_A(q)_{A+c})=\sum_{m=0}^{A}q^{m(m+c)}/((q)_m(q)_{m+c}(q)_{A-m})$, and as $A\to\infty$
(coefficientwise in $\Z[[q]]$) the left side tends to $(q)_\infty^{-2}$ and the right side to
$(q)_\infty^{-1}\sum_mq^{m(m+c)}/((q)_m(q)_{m+c})$. Use it with $c=R+a+j$ for each $j$ and substitute $m\to m-j$:
\[
\frac1{(q)_\infty}\sum_j\qbin ajq^{j^2+Rj}=\sum_{m\ge0}\frac{q^{m^2+Rm}}{(q)_{m+R+a}}\sum_{j=0}^{\min(a,m)}\qbin aj\frac{q^{(m-j)(a-j)}}{(q)_{m-j}},
\]
because $j^2+Rj+(m-j)(m+R+a)=m^2+Rm+(m-j)(a-j)$. By Lemma~\ref{lem:durfee},
$\sum_jq^{(m-j)(a-j)}/((q)_j(q)_{a-j}(q)_{m-j})=1/((q)_a(q)_m)$, so the inner sum is $1/(q)_m$.
\end{proof}

\begin{proposition}\label{prop:T}
For integers $k\ge1$ and $a\ge0$,
\[
T_k(a)=\frac1{(q)_\infty}\sum_{r\ge0}(-1)^rq^{kar+(k+1)r^2+r(r-1)/2}\,\frac{(1-q^{a+2r})(q)_{a+r-1}}{(q)_r}\sum_{j=0}^{a}\qbin ajq^{j^2+2rj},
\]
where for $a=r=0$ the factor $(1-q^{a+2r})(q)_{a+r-1}/(q)_r$ is read as $1$.
\end{proposition}
\begin{proof}
With $x=q^a$ and $N=N_1$, the left side of \eqref{eq:bailey} multiplied by $q^{N^2}$ and summed over $N$ is
$T_k(a)$ (for $k=1$, $T_1(a)=\sum_Nq^{N^2}/(q)_N$). Interchanging the summable sums and writing $N=m+r$,
\[
T_k(a)=\sum_{r\ge0}q^{kar+kr^2}\alpha_r\sum_{m\ge0}\frac{q^{(m+r)^2}}{(q)_m(q^{a+1};q)_{m+2r}} .
\]
Since $(q^{a+1})_{m+2r}=(q)_{m+2r+a}/(q)_a$, Lemma~\ref{lem:cauchy} with $R=2r$ gives for the inner sum
$q^{r^2}(q)_a(q)_\infty^{-1}\sum_j\qbin ajq^{j^2+2rj}$. Finally
$\alpha_r(q)_a=(-1)^rq^{r(r-1)/2}(1-q^{a+2r})(q)_{a+r-1}/(q)_r$, because $(q^a)_r/(1-q^a)=(q)_{a+r-1}/(q)_a$ for
$a\ge1$, and by the convention above for $a=0$.
\end{proof}

Insert Proposition~\ref{prop:T} in the left side of Theorem~\ref{thm:family}, which is
$\sum_aq^{ka^2/4}T_k(a)/(q)_a$. The triple series over $(a,r,j)$ is summable (the exponent is
$\ge\frac k4(a+2r)^2$), so we may collect the terms with $n:=a+2r$ fixed. Since
$\frac k4a^2+kar+(k+1)r^2+\frac{r(r-1)}2+j^2+2rj=\frac k4n^2+\frac{r(r-1)}2+(r+j)^2$ and
$\qbin aj/(q)_a=1/((q)_j(q)_{a-j})$, we get
\begin{equation}\label{eq:cn}
(q)_\infty\cdot(\text{left side})=\sum_{n\ge0}q^{kn^2/4}c_n,\qquad c_0=1,\end{equation}
\begin{equation}\label{eq:cn2}
c_n=(1-q^n)\!\!\sum_{\substack{r,j\ge0\\2r+j\le n}}\!\!\frac{(-1)^rq^{r(r-1)/2+(r+j)^2}(q)_{n-r-1}}{(q)_r(q)_j(q)_{n-2r-j}}\quad(n\ge1).
\end{equation}
The polynomials $c_n$ do not depend on $k$.

\begin{proposition}\label{prop:cn}
$c_n=1+q^{n^2}$ for all $n\ge1$.
\end{proposition}
\begin{proof}
Put $u=r+j$. Then
\[
\frac{c_n}{1-q^n}=\sum_{u=0}^nq^{u^2}\Sigma_u,\qquad
\Sigma_u:=\sum_{r=0}^{\min(u,n-u)}\frac{(-1)^rq^{r(r-1)/2}(q)_{n-r-1}}{(q)_r(q)_{u-r}(q)_{n-u-r}} .
\]
For $u=0$ and $u=n$ only $r=0$ occurs, and $\Sigma_0=\Sigma_n=(q)_{n-1}/(q)_n=1/(1-q^n)$. For $1\le u\le n-1$ use
$(q)_{n-r-1}=(q)_{n-1}/(q^{n-r};q)_r$ with $(q^{n-r};q)_r=(-1)^rq^{nr-r(r+1)/2}(q^{1-n};q)_r$,
$1/(q)_{u-r}=(-1)^rq^{ur-r(r-1)/2}(q^{-u};q)_r/(q)_u$ and $1/(q)_{n-u-r}=(-1)^rq^{(n-u)r-r(r-1)/2}(q^{u-n};q)_r/(q)_{n-u}$;
the powers of $q$ combine to $q^r$ and the signs to $+1$:
\[
\Sigma_u=\frac{(q)_{n-1}}{(q)_u(q)_{n-u}}\sum_{r\ge0}\frac{(q^{-u};q)_r(q^{u-n};q)_r}{(q;q)_r(q^{1-n};q)_r}q^r
=\frac{(q)_{n-1}}{(q)_u(q)_{n-u}}\,{}_2\phi_1(q^{-u},q^{u-n};q^{1-n};q,q)
\]
(the denominators $(q^{1-n};q)_r$, $r\le u\le n-1$, do not vanish). By the $q$-Chu--Vandermonde sum
${}_2\phi_1(q^{-u},b;c;q,q)=(c/b;q)_ub^u/(c;q)_u$~\cite[(1.5.3)]{GR} (an identity of rational functions of
$b,c$, valid when $(c;q)_u\ne0$) with $b=q^{u-n}$, $c=q^{1-n}$ we have $c/b=q^{1-u}$ and
$(q^{1-u};q)_u=0$ for $u\ge1$. Hence $\Sigma_u=0$ for $1\le u\le n-1$ and $c_n=(1-q^n)(1+q^{n^2})/(1-q^n)$.
\end{proof}

\begin{proof}[Proof of Theorems~\ref{thm:family} and~\ref{thm:E7}]
By \eqref{eq:cn} and Proposition~\ref{prop:cn},
$(q)_\infty\cdot(\text{left side})=1+\sum_{n\ge1}(q^{kn^2/4}+q^{kn^2/4+n^2})$, which is
Theorem~\ref{thm:family}. (The coefficient $k/4$ of $a^2$ is exactly what makes the cross term in $nr$
cancel.) For $k=3$, \eqref{eq:E7red} gives the first formula of Theorem~\ref{thm:E7}; the second follows
since $\sum_{\Z}q^{3k^2/4}+\sum_\Z q^{7k^2/4}=2+2\sum_{k\ge1}(\cdots)$.

\emph{Modularity.} Put $\theta(z)=\sum_{k\in\Z}e^{2\pi ik^2z}$, a modular form of weight $\frac12$ on
$\Gamma_0(4)$ (with the theta multiplier), and $\eta(z)=e^{2\pi iz/24}\prod(1-e^{2\pi inz})$. With
$q=e^{2\pi i\tau}$ and $c_{\rm eff}=1$, $q^{-1/24}F(q)=(\theta(3\tau/4)+\theta(7\tau/4))/(2\eta(\tau))$; in the
variable $z=\tau/24$ this is $(\theta(18z)+\theta(42z))/(2\eta(24z))$. Here $\theta(Nz)$ is a modular form of
weight $\frac12$ on $\Gamma_0(4N)$ with the theta multiplier times the quadratic character $(N/\cdot)$, and
$\eta(24z)=\sum_{n\ge1}\chi_{12}(n)q^{n^2}$ (Euler's pentagonal theorem, $q=e^{2\pi iz}$) is a modular form of
weight $\frac12$ on $\Gamma_0(576)$ with the theta multiplier times $\chi_{12}$ (see~\cite{Shimura}
and~\cite[Ch.~1]{Ono}). The three characters are trivial on $\Gamma_1(4032)$,
$4032=\mathrm{lcm}(72,168,576)$, so $\theta(18z)+\theta(42z)$ and $\eta(24z)$ are weight-$\frac12$ forms on
$\Gamma_1(4032)$ with the same (theta) multiplier, and their quotient is a weakly holomorphic modular function
on $\Gamma_1(4032)$; in the variable $\tau=24z$ it is invariant under $\Gamma(4032)$.
\end{proof}

\emph{The parity refinement.} The diagonal of $C(E_7)^{-1}$ is $(3/2,4,15/2,12,6,2,7/2)$ and its
off-diagonal half-integers are exactly between two of the nodes $1,3,7$; so
$\frac12n^{\T}C(E_7)^{-1}n\equiv\frac34(n_1+n_3+n_7)^2\pmod1$, and the exponents lie in $\Z$ or in
$\frac34+\Z$ according to the parity of $n_1+n_3+n_7$. Since $3k^2/4,7k^2/4\in\Z$ iff $k$ is even, separating
the two classes on both sides gives the refinements stated after Theorem~\ref{thm:E7}.

\begin{remark}[the cases $k=1,2$ of Theorem~\ref{thm:family}]
For $k=1$, $T_1(a)$ does not depend on $a$, and the theorem reads
$(\sum_aq^{a^2/4}/(q)_a)\,G(q)=(q)_\infty^{-1}(\sum_{m\ge0}q^{m^2/4}+\sum_{m\ge1}q^{5m^2/4})$. For $k=2$
the matrix $\left(\begin{smallmatrix}1&0&1\\0&2&2\\1&2&4\end{smallmatrix}\right)$ is one fusion away from
$C(T_4)^{-1}=(\min(i,j))_{i,j\le4}$, so the left side is the $(T_1,T_4)$ sum
$\sum q^{\frac12(M_1^2+\dots+M_4^2)}/\prod(q)_{m_i}$: this is Melzer's identity, proved
in~\cite{BIS,Warnaar03} as recalled in §\ref{ssec:novelty}. For $k\ge3$ these sums differ from Melzer's: the
terminal form of $C(T_{2k})^{-1}$ under defusion has $A_{aa}=1$ and $a$ coupled to $N_k$ only, whereas here
$A_{aa}=k/2$ and $a$ is coupled to $N_2+\dots+N_k$; the two families meet only at $k=2$.
\end{remark}

\section{Further remarks and open problems}\label{sec:open}

\begin{enumerate}[leftmargin=2em]
\item \emph{Known results by the same method} (sketches, not used above). For $r\ge3$, every chain node of
$2C(D_r)^{-1}$ is successively the fusion of the two spinor nodes; removing them leaves the rank-two matrix
$\left(\begin{smallmatrix}r/2&1-r/2\\1-r/2&r/2\end{smallmatrix}\right)$, and Cauchy's identity gives
$\sum_nq^{n^{\T}C(D_r)^{-1}n}/\prod(q)_{n_i}=(q)_\infty^{-1}\sum_{k\in\Z}q^{rk^2/4}$ (the
Kedem--Klassen--McCoy--Melzer identity~\cite{KKMM1}; see~\cite[Theorem~1.2]{WW-Dk} and the references
there; checked for $r\le8$).
Similarly $C(T_{2r+1})^{-1}$ defuses to $(1)\oplus2C(T_r)^{-1}$ (checked for $r\le3$), which gives the
$(T_1,T_{2r+1})$ identity of Warnaar~\cite[Theorem~4.2]{Warnaar03} (Theorem~4.5 in the numbering quoted
in~\cite{SW}) from Andrews--Gordon.
\item \emph{The dual pairs $(E_n,T_1)$.} For the Nahm sums of $C(E_n)$ itself fusion does not apply: reading
Lemma~\ref{lem:fusion} backwards would need three nodes $a,b,c$ of a Cartan matrix with
$A_{ac}=A_{aa}+A_{ab}-1$ and $A_{cc}=A_{aa}+A_{bb}+2(A_{ab}-1)$, i.e.\ $A_{ac}=1>0$ if $A_{ab}=0$ and
$A_{cc}=0\ne2$ if $A_{ab}=-1$. Sun and Wang~\cite[(3.21)]{SW} identify the $(E_8,T_1)$ sum numerically with the vacuum
character of the Hecke image $T_{10}M_{\rm eff}(11,2)$; Yoshida~\cite{Yoshida} gives four companions. We
checked all five against the Hecke images computed from the formula of Harvey and Wu~\cite{HW} modulo
$q^{26}$; they remain unproved. For $(E_7,T_1)$ we found, by analogy with Theorem~\ref{thm:E7},
\begin{equation}\label{eq:E7T1}
\sum_{n\in\N^7}\frac{q^{\frac12n^{\T}C(E_7)n}}{\prod(q)_{n_i}}\overset{?}{=}\frac{\Theta_{E_6}(q)+\Theta_{A_6}(q)}{2(q;q)_\infty^6},
\qquad\Theta_L(q)=\sum_{x\in L}q^{|x|^2/2},
\end{equation}
which holds modulo $q^{40}$ (the analogy: the even part of Theorem~\ref{thm:E7} is
$(\Theta_{\langle6\rangle}+\Theta_{\langle14\rangle})/(2(q)_\infty)$, and Theorem~\ref{thm:E7} itself is
$(\Theta_{\langle6\rangle}+\Theta_{\langle14\rangle})(q^{1/4})/(2(q)_\infty)$, where $\langle6\rangle$ and $\langle14\rangle$
are the rank-one lattices orthogonal to $E_6$ and to $A_6$ in $E_7$). We state \eqref{eq:E7T1} as a conjecture. For $(E_6,T_1)$ we found no formula.
\item \emph{Rank of the reduced forms.} The reductions are to ranks $4$ ($E_8$), $3$ ($E_6$) and $4$ ($E_7$),
and Theorem~\ref{thm:family} is a one-parameter family containing the $E_7$ case. It would be interesting
to know which other Nahm sums $C(X)\otimes C(Y)^{-1}$ reduce by fusion to sums of known type.
\end{enumerate}

\section{The Lean formalization}\label{sec:lean}

The Lean~4 project \fpath{RRFusion}, in the public repository at
\repo\fpath{02-exceptional-nahm-sums/lean/} (pins: Lean \texttt{leanprover/lean4:v4.33.0-rc2}, Mathlib
commit \texttt{51e6992efd06\allowbreak126df61a496bebf8\allowbreak f49482a4e129}), machine-checks the finite and algebraic
ingredients of §§\ref{sec:lemmas}--\ref{sec:E6}. It uses no \texttt{sorry}, no additional axiom and no
\texttt{native\_decide}; \verb|#print axioms| gives \texttt{[propext, Classical.choice, Quot.sound]} for
the sixteen main theorems below (the auxiliary ones were not included in that run). Matrices of Nahm data are functions $\iota\to\iota\to\Q$ on a finite index type
$\iota$; fusing two species enlarges the index type to \texttt{Option}~$\iota$, the new species being
\texttt{none}.
Theorems (namespace \texttt{RRFusion}):
\begin{itemize}[leftmargin=1.5em]
\item \texttt{C8\_mul\_A8}, \texttt{A8\_mul\_C8} (by \texttt{decide}, i.e.\ kernel evaluation over $\Z$): the
printed $8\times8$ matrix is a two-sided inverse of the $E_8$ Cartan matrix; \texttt{C6\_mul\_A6},
\texttt{A6\_mul\_C6} (entrywise, by \texttt{fin\_cases}, \texttt{simp}, \texttt{norm\_num} over $\Q$): the same for
$E_6$. Auxiliary: \texttt{C8\_eq\_diagram}, \texttt{C6\_eq\_diagram} (the matrices are those of the diagrams
of §\ref{sec:intro}), \texttt{C8\_eq\_mathlib}, \texttt{C6\_eq\_mathlib} (they are Mathlib's
\texttt{CartanMatrix.E}$_8$, \texttt{CartanMatrix.E}$_6$ after an explicit relabelling), \texttt{A8\_symm},
\texttt{A6\_symm}.
\item \texttt{Q\_fuse}: identity \eqref{eq:Qfuse}, for an arbitrary finite index type, symmetric $A$, any
$B$, $a\ne b$, $m\in\Q^\iota$, $c\in\Q$ (a general proof); \texttt{fuse\_symm}: fusion preserves symmetry.
\item \texttt{E8\_chain}, \texttt{E8\_chain\_linear}, \texttt{nodeE8\_bijective}: the four fusions of
§\ref{sec:E8} applied to $(2\min(i,j))$ give, under a bijective relabelling of the nodes, exactly the printed
$C(E_8)^{-1}$ and the vector $B(b)$ of Theorem~\ref{thm:E8} for every $b\in\Q^4$; \texttt{E6\_chain},
\texttt{E6\_chain\_linear}, \texttt{nodeE6\_bijective}: the same for $E_6$.
\item \texttt{durfee\_gauss}: in any commutative ring,
$\sum_{c=0}^{\min(a,b)}q^{(a-c)(b-c)}\qbin ac\qbin bc(q)_c=1$; \texttt{durfee\_inv}, \texttt{durfee\_field},
\texttt{durfee\_ratFunc}: identity \eqref{eq:durfee} in a field with $q^n\ne1$ ($n>0$), in particular in
$\Q(q)$. (Proved by induction on the $q$-Pascal recurrence, not by the generating-function proof of
Lemma~\ref{lem:durfee}.)
\end{itemize}
\emph{Not formalized} (as of this version): Lemma~\ref{lem:fusion} for infinite sums (summability,
admissibility, the regrouping); the Andrews--Gordon, Rogers--Ramanujan and Cao--Rosengren--Wang identities;
the power-series identities of Theorems~\ref{thm:E8}--\ref{thm:family}; the $E_7$ chain and
§§\ref{sec:E7}--\ref{sec:family}; modularity. A separate effort is formalizing the full identities. To
rebuild: install \texttt{elan}, then in the project folder run \texttt{lake exe cache get} and
\texttt{lake build}.

\section{Numerical verification}\label{sec:num}
None of the following is used in the proofs. Paths are relative to
\repo\fpath{02-exceptional-nahm-sums/}.
\begin{itemize}[leftmargin=1.5em]
\item \fpath{code/certify_E8_E6.py} (log \fpath{certificates/certify_E8_E6.txt}): the fusion chains and
inverse matrices exactly; Corollary~\ref{cor:E8} for $i=1,\dots,5$ modulo $q^{300}$; Corollary~\ref{cor:E6}
(all six), \eqref{eq:Ex10} and the character form of \eqref{eq:SW319} modulo $q^{150}$.
\fpath{code/verify_E8_E6.py} (log \fpath{certificates/verify_E8_E6_q500.txt}): \eqref{eq:SW320} modulo
$q^{500}$ (72\,395 lattice points) and \eqref{eq:SW319} modulo $q^{300}$. \fpath{code/testE8y.gp}
(PARI/GP, log \fpath{certificates/testE8y.txt}): Theorem~\ref{thm:E8} for ten integer vectors $b$,
including negative entries, modulo $q^{70}$.
\item \fpath{code/certify_E7.py} (log \fpath{certificates/certify_E7.txt}): the $E_7$ chain exactly;
Proposition~\ref{prop:cn} exactly for $n\le30$; Proposition~\ref{prop:T} ($k=3$) for $a\le5$ modulo
$q^{120}$; Theorem~\ref{thm:E7} as a series in $q^{1/4}$ modulo $q^{300}$. \fpath{code/test_family_k.py}
(log \fpath{certificates/test_family_k.txt}): Theorem~\ref{thm:family} for $k\le6$ modulo $q^{80}$.
\item \fpath{code/sw_E7T1.py} (log \fpath{certificates/sw_E7T1.txt}): \eqref{eq:E7T1} modulo $q^{40}$;
\fpath{code/sw_yoshida.py} (log \fpath{certificates/sw_yoshida.txt}): the $(E_8,T_1)$ comparison modulo $q^{26}$; \fpath{code/sw_audit.py}
(log \fpath{certificates/sw_audit.txt}): Appendix~\ref{app:SW}.
\item Independent referee computations (reports \fpath{REFEREE-E8-E6.md}, \fpath{REFEREE-2-E8-E6.md},
\fpath{REFEREE-E7-KRTW.md}, \fpath{REFEREE-2-E7.md}; code written from the statements, not archived in
the repository; the reports describe the methods):
Corollary~\ref{cor:E8} modulo $q^{400}$; Theorem~\ref{thm:E8} for thirteen rational vectors $b$ modulo
$q^{50}$; Corollary~\ref{cor:E6} to $q^{147}$; Theorem~\ref{thm:E7} by direct enumeration of $\N^7$ through
$q^{200}$; Theorem~\ref{thm:family} for $k\le8$ through $q^{150}$; positive definiteness of all
intermediate matrices.
\end{itemize}

\section*{Acknowledgements}
The results, proofs, computer code and Lean formalization in this paper were produced by AI (Anthropic's
Claude models) directed by the author. Their correctness was checked by independent AI referees
(separate instances instructed to find errors, each writing its own code; four reports, included in the
repository) and, for the formalized statements, by the Lean kernel. No human expert has refereed this
work yet; corrections are welcome. We thank Kaiwen Sun and Haowu Wang for their extensive and carefully
tested tables, which made this work possible.

\appendix
\section{Remarks on the tables of Sun and Wang}\label{app:SW}

While testing the identities of~\cite{SW} we compared all displayed formulas (1.8)--(1.12), (2.1)--(2.8),
(3.2), (3.5)--(3.41) and the central charges of Table~1 numerically with the Nahm sums (137 series
comparisons, to orders between about $q^{50}$ and $q^{500}$; \fpath{certificates/sw_audit.txt}). All of them
agree except for three small misprints, which we record for the convenience of readers; none affects the
claims of~\cite{SW}.
\begin{enumerate}[leftmargin=2em]
\item In (3.27) and (3.28) (the pair $(T_2,T_1)$) the Ramond weights $1/32$ and $5/32$ appear to be
interchanged, and the first identity needs a factor $2$: the Nahm sums satisfy
$f_3=2\chi_{R,5/32}\chi^F_{R,1/16}$ and $f_4=\chi_{R,1/32}\chi^F_{R,1/16}$ (checked modulo $q^{300}$); as printed,
the leading terms do not match ($f_3=2q^{1/6}+\cdots$).
\item In the sentence before (3.30), ``the Nahm sum associated with $(T_1,A_3)$ is identical to $f_1$''
should read $f_2$; formula (3.30) itself is correct.
\item In (3.6) the same line is printed for $r=2n$ and $r=2n+1$. This is harmless: the formula
$\chi^{\rm NS}_0-\chi^{\rm NS}_{\lfloor r/2\rfloor+1}$ holds for both parities (checked for $r=2,\dots,9$); the Kac
labels differ between the two cases, the conformal weights do not.
\end{enumerate}
The formulas with a coefficient $2$ in front of a Ramond character ((1.8)--(1.11), (3.30)) hold when
$\chi_R$ is normalised with leading coefficient $1$; we mention this convention because it is not stated
explicitly. One further failure in our comparison is the value $C=19/24$ in Zagier's original list, which
Sun and Wang themselves correct to $19/96$.

\begin{thebibliography}{99}

\bibitem{Andrews74} G.~E. Andrews, \emph{An analytic generalization of the Rogers--Ramanujan identities
for odd moduli}, Proc. Nat. Acad. Sci. USA \textbf{71} (1974), 4082--4085.

\bibitem{Andrews84} G.~E. Andrews, \emph{Multiple series Rogers--Ramanujan type identities}, Pacific J.
Math. \textbf{114} (1984), 267--283.

\bibitem{AndrewsBook} G.~E. Andrews, \emph{The Theory of Partitions}, Encyclopedia of Mathematics and its
Applications~2, Addison-Wesley, 1976; reprinted Cambridge University Press, 1998.

\bibitem{BIS} D.~Bressoud, M.~Ismail and D.~Stanton, \emph{Change of base in Bailey pairs}, Ramanujan J.
\textbf{4} (2000), 435--453; arXiv:math/9909053.

\bibitem{CRW} Z.~Cao, H.~Rosengren and L.~Wang, \emph{On some double Nahm sums of Zagier}, J. Combin.
Theory Ser. A \textbf{202} (2024), Paper No.~105819; arXiv:2303.01333.

\bibitem{CW-lift3} Z.~Cao and L.~Wang, \emph{Some new modular rank three Nahm sums from a lift-dual
operation}, arXiv:2412.15767 (2024).

\bibitem{CW-lift4} Z.~Cao and L.~Wang, \emph{Some new modular rank four Nahm sums as lift-dual of rank
three examples}, arXiv:2508.12468 (2025).

\bibitem{CW3} Z.~Cao and L.~Wang, \emph{Some new modular Nahm sums of ranks 3 and 4}, arXiv:2606.13590
(2026).

\bibitem{CF} I.~Cherednik and B.~Feigin, \emph{Rogers--Ramanujan type identities and Nil-DAHA}, Adv.
Math. \textbf{248} (2013); arXiv:1209.1978.

\bibitem{Lean4} L.~de~Moura and S.~Ullrich, \emph{The Lean~4 theorem prover and programming language},
in: Automated Deduction -- CADE 28, Lecture Notes in Comput. Sci. \textbf{12699}, Springer, 2021,
625--635.

\bibitem{EKL} T.~Ekholm, P.~Kucharski and P.~Longhi, \emph{Multi-cover skeins, quivers, and 3d
$\mathcal N=2$ dualities}, J. High Energy Phys. 2020, no.~2, 018; arXiv:1910.06193.

\bibitem{GR} G.~Gasper and M.~Rahman, \emph{Basic Hypergeometric Series}, 2nd ed., Encyclopedia of
Mathematics and its Applications~96, Cambridge University Press, 2004.

\bibitem{Forty} B.~C. Berndt, G.~Choi, Y.-S. Choi, H.~Hahn, B.~P. Yeap, A.~J. Yee, H.~Yesilyurt and
J.~Yi, \emph{Ramanujan's forty identities for the Rogers--Ramanujan functions}, Mem. Amer. Math. Soc.
\textbf{188} (2007), no.~880.

\bibitem{Gordon} B.~Gordon, \emph{A combinatorial generalization of the Rogers--Ramanujan identities},
Amer. J. Math. \textbf{83} (1961), 393--399.

\bibitem{HW} J.~A. Harvey and Y.~Wu, \emph{Hecke relations in rational conformal field theory},
arXiv:1804.06860 (2018).

\bibitem{KKMM1} R.~Kedem, T.~R. Klassen, B.~M. McCoy and E.~Melzer, \emph{Fermionic quasi-particle
representations for characters of $(G^{(1)})_1\times(G^{(1)})_1/(G^{(1)})_2$}, Phys. Lett. B \textbf{304}
(1993), 263--270; arXiv:hep-th/9211102.

\bibitem{KKMM2} R.~Kedem, T.~R. Klassen, B.~M. McCoy and E.~Melzer, \emph{Fermionic sum representations for
conformal field theory characters}, Phys. Lett. B \textbf{307} (1993), 68--76; arXiv:hep-th/9301046.

\bibitem{KLNS} P.~Kucharski, H.~Larragu\'ivel, D.~Noshchenko and P.~Su\l kowski, \emph{Unlinking symmetric
quivers}, arXiv:2312.14905 (2023).

\bibitem{Mathlib} The mathlib Community, \emph{The Lean mathematical library}, in: Proceedings of the 9th
ACM SIGPLAN International Conference on Certified Programs and Proofs (CPP 2020), ACM, 2020, 367--381.

\bibitem{Melzer} E.~Melzer, \emph{Supersymmetric analogs of the Gordon--Andrews identities, and related
TBA systems}, arXiv:hep-th/9412154 (1994).

\bibitem{Nahm} W.~Nahm, \emph{Conformal field theory and torsion elements of the Bloch group}, in:
Frontiers in Number Theory, Physics, and Geometry~II, Springer, Berlin, 2007, 67--132;
arXiv:hep-th/0404120.

\bibitem{Ono} K.~Ono, \emph{The Web of Modularity: Arithmetic of the Coefficients of Modular Forms and
$q$-series}, CBMS Regional Conference Series in Mathematics~102, Amer. Math. Soc., 2004.

\bibitem{RC} A.~Rocha-Caridi, \emph{Vacuum vector representations of the Virasoro algebra}, in: Vertex
Operators in Mathematics and Physics, Math. Sci. Res. Inst. Publ.~3, Springer, New York, 1985, 451--473.

\bibitem{Rogers} L.~J. Rogers, \emph{Second memoir on the expansion of certain infinite products}, Proc.
London Math. Soc. \textbf{25} (1894), 318--343.

\bibitem{Shimura} G.~Shimura, \emph{On modular forms of half integral weight}, Ann. of Math. (2)
\textbf{97} (1973), 440--481.

\bibitem{SW} K.~Sun and H.~Wang, \emph{Dynkin diagrams, generalized Nahm sums and 2d CFTs},
arXiv:2604.00847 (2026).

\bibitem{VZ} M.~Vlasenko and S.~Zwegers, \emph{Nahm's conjecture: asymptotic computations and
counterexamples}, Commun. Number Theory Phys. \textbf{5} (2011), no.~3, 617--642; arXiv:1104.4008.

\bibitem{Wang22} L.~Wang, \emph{Identities on Zagier's rank two examples for Nahm's problem},
arXiv:2210.10748 (2022).

\bibitem{WW-Dk} L.~Wang and S.~Wang, \emph{Nahm sum identities for Cartan matrices of type $D_k$},
arXiv:2512.07790 (2025).

\bibitem{WW-Dr} L.~Wang and S.~Wang, \emph{Modular Nahm sums for the inverse Cartan matrix of type $D_r$},
arXiv:2607.08606 (2026).

\bibitem{Warnaar50} S.~O. Warnaar, \emph{50 years of Bailey's lemma}, in: Algebraic Combinatorics and
Applications (A.~Betten et al., eds.), Springer, Berlin, 2001, 333--347; arXiv:0910.2062.

\bibitem{Warnaar03} S.~O. Warnaar, \emph{The generalized Borwein conjecture. II. Refined $q$-trinomial
coefficients}, Discrete Math. \textbf{272} (2003), 215--258; arXiv:math/0110307 (theorem numbers refer to
the arXiv version).

\bibitem{WP} S.~O. Warnaar and P.~A. Pearce, \emph{Exceptional structure of the dilute $A_3$ model: $E_8$
and $E_7$ Rogers--Ramanujan identities}, J. Phys. A \textbf{27} (1994), L891--L898; arXiv:hep-th/9408136.

\bibitem{Yoshida} Y.~Yoshida, \emph{3d $\mathcal N=4$ rank-zero mirror symmetry, TQFT interfaces, and
Zagier duality of Nahm sums}, arXiv:2608.08780 (2026).

\bibitem{Zagier} D.~Zagier, \emph{The dilogarithm function}, in: Frontiers in Number Theory, Physics, and
Geometry~II, Springer, Berlin, 2007, 3--65.

\end{thebibliography}
\end{document}
